\begin{table}[htbp]
\centering
\caption{AUROC and AUPRC with 95% stratified bootstrap confidence intervals, per model x calibration method (held-out test set, n=2158, full 21,388-record dataset).}
\label{table_discrimination_auroc_auprc_ci}
\begin{tabular}{llrll}
\toprule
model & calibration\_method & n\_test & auroc & auprc \\
\midrule
lightgbm & dwac & 2158 & 0.8464 (0.8304–0.8618) & 0.8932 (0.8809–0.9043) \\
lightgbm & isotonic & 2158 & 0.8457 (0.8291–0.8610) & 0.8822 (0.8691–0.8941) \\
lightgbm & platt & 2158 & 0.8464 (0.8305–0.8618) & 0.8932 (0.8809–0.9043) \\
lightgbm & uncalibrated & 2158 & 0.8464 (0.8305–0.8618) & 0.8932 (0.8809–0.9043) \\
logistic\_regression & dwac & 2158 & 0.7953 (0.7762–0.8146) & 0.8480 (0.8311–0.8636) \\
logistic\_regression & isotonic & 2158 & 0.7933 (0.7744–0.8121) & 0.8333 (0.8164–0.8496) \\
logistic\_regression & platt & 2158 & 0.7954 (0.7761–0.8142) & 0.8480 (0.8311–0.8636) \\
logistic\_regression & uncalibrated & 2158 & 0.7954 (0.7761–0.8142) & 0.8480 (0.8311–0.8636) \\
\bottomrule
\end{tabular}

\end{table}


\begin{quote}\small
\textit{Abbreviations:} AUROC = Area Under the Receiver Operating Characteristic curve; AUPRC = Area Under the Precision-Recall Curve; CI = Confidence Interval.
\textit{Notes:} 1000 stratified bootstrap resamples (resampled within each outcome class, preserving observed class counts), seed=42. AUROC/AUPRC are numerically identical across Platt/DWAC/uncalibrated for a given model where the calibration transform preserves rank order; isotonic regression can introduce small AUROC/AUPRC differences due to tied outputs in flat regions of the fitted step function.
\end{quote}